%!TEX program = xelatex
\PassOptionsToPackage{full}{textcomp}
\documentclass[lang=en,11pt,a4paper,citestyle=authoryear]{elegantpaper}

\title{HW6 - STAT 709}
\author{Wells Guan}

\usepackage{array,url,stix}
\usepackage{tikz}
\usepackage{tikz-cd}
\usepackage{enumitem}
\usepackage{amsmath}
\usepackage{amssymb}
\usepackage{listings}
\usepackage{subfigure}

\lstset{
    frame=single,
    columns=fullflexible,
    keepspaces=true,
    showstringspaces=false,
    breaklines=true,
    breakatwhitespace=false,
    basicstyle=\ttfamily\footnotesize,
    numbers=left,
    numberstyle=\tiny,
    numbersep=6pt
}

\newcommand{\Z}{\mathbb{Z}}
\newcommand{\R}{\mathbb{R}}
\newcommand{\N}{\mathbb{N}}
\newcommand{\B}{\mathcal{B}}
\newcommand{\E}{\mathbb{E}}
\renewcommand{\P}{\mathbb{P}}
\newcommand{\norm}[1]{\left\lVert#1\right\rVert}

\begin{document}
\maketitle

\noindent All exercises are required unless marked optional.

% Exercise statements below were transcribed from the supplied screenshots.
% Exercise 5.3.33(b) incorporates the correction in the assignment PDF.
% Write solutions after the Sol. headings.

\subsection*{Problem 1. Modes of convergence}
Jun Shao, \textit{Mathematical Statistics}, second edition, Chapter 1:
\begin{itemize}[leftmargin=*]
\item Required: Exercises 127, 135, and 140.
\item Optional: Exercises 137 and 152.
\end{itemize}

\subsubsection*{Exercise 1.127}
Let $X,X_1,X_2,\ldots$, $Y_1,Y_2,\ldots$, and $Z_1,Z_2,\ldots$
be random variables. Prove the following statements.
\begin{enumerate}[label=(\alph*),leftmargin=*]
\item If $X_n\xrightarrow{d}X$, then $X_n=O_p(1)$.
\item If $X_n=O_p(Z_n)$ and $\P(Y_n=0)=0$, then
$X_nY_n=O_p(Y_nZ_n)$.
\item If $X_n=O_p(Z_n)$ and $Y_n=O_p(Z_n)$, then
$X_n+Y_n=O_p(Z_n)$.
\item If $\E|X_n|=O(a_n)$, then $X_n=O_p(a_n)$, where
$a_n\in(0,\infty)$.
\item If $X_n\xrightarrow{\mathrm{a.s.}}X$, then
$\sup_n|X_n|=O_p(1)$.
\end{enumerate}

\par\vspace{0.5em}\noindent\textbf{Sol.}\par
Recall that $X_n=O_p(Z_n)$ means that, for every $\epsilon>0$,
there are $M<\infty$ and $N$ such that
$\P(|X_n|>M|Z_n|)<\epsilon$ for all $n\ge N$.
This formulation avoids dividing by a possibly zero random scale.
\begin{enumerate}[label=(\alph*),leftmargin=*]
\item For any $\epsilon>0$, choose $M>0$ so that
\[
\P(|X|>M)<\epsilon/2,
\qquad F\text{ is continuous at }M\text{ and }-M,
\]
where $F(x)=\P(X\le x)$ and $F_n(x)=\P(X_n\le x)$.
Such an $M$ exists because $X$ is finite almost surely and a CDF
has at most countably many discontinuities.
Convergence in distribution gives an integer $N$ such that
\[
|F_n(M)-F(M)|+|F_n(-M)-F(-M)|<\epsilon/2,
\qquad n\ge N.
\]
Consequently,
\[
\begin{aligned}
\P(|X_n|>M)
&\le 1-F_n(M)+F_n(-M)\\
&\le 1-F(M)+F(-M)+\epsilon/2\\
&=\P(|X|>M)+\epsilon/2<\epsilon.
\end{aligned}
\]
The last equality uses $\P(X=-M)=0$, while the first line is an
inequality because $F_n(-M)$ also counts a possible atom at $-M$.
Thus $X_n=O_p(1)$.

\item Since $\P(Y_n=0)=0$, multiplication by $|Y_n|$ preserves the
relevant comparison almost surely, so for every $M>0$,
\[
\P(|X_nY_n|>M|Z_nY_n|)
=\P(|X_n|>M|Z_n|).
\]
The right side is eventually less than any prescribed $\epsilon>0$
for a suitable $M$, which proves $X_nY_n=O_p(Y_nZ_n)$.

\item For any $\epsilon>0$, choose $M_1,M_2$ so that, for all
sufficiently large $n$,
\[
\P(|X_n|>M_1|Z_n|)<\epsilon/2,
\qquad
\P(|Y_n|>M_2|Z_n|)<\epsilon/2.
\]
By the triangle inequality and a union bound,
\[
\begin{aligned}
\P(|X_n+Y_n|>(M_1+M_2)|Z_n|)
&\le\P(|X_n|>M_1|Z_n|)\\
&\quad+\P(|Y_n|>M_2|Z_n|)<\epsilon.
\end{aligned}
\]
Therefore $X_n+Y_n=O_p(Z_n)$; no independence is needed.

\item By $\E|X_n|=O(a_n)$, there are $K<\infty$ and $N_0$ such that
$\E|X_n|\le Ka_n$ for $n\ge N_0$. Markov's inequality yields
\[
\P(|X_n|>Ma_n)\le\frac{\E|X_n|}{Ma_n}\le\frac KM.
\]
Choosing $M>K/\epsilon$ proves $X_n=O_p(a_n)$.

\item On the event where $X_n$ converges to the finite random variable
$X$, the numerical sequence $(X_n)$ is bounded. Thus
\[
S:=\sup_{n\ge1}|X_n|<\infty\qquad\text{a.s.}
\]
The supremum is measurable because it is countable. Therefore
$\P(S>M)\to0$ as $M\to\infty$, so the constant sequence with value
$S$ is $O_p(1)$. In fact, the same estimate also shows
$\sup_n\P(|X_n|>M)\le\P(S>M)\to0$.
\end{enumerate}
\vspace{2em}

\subsubsection*{Exercise 1.135}
Let $X,X_1,X_2,\ldots$ be random $k$-vectors and
$Y,Y_1,Y_2,\ldots$ be random $l$-vectors. Suppose that
$X_n\xrightarrow{d}X$, $Y_n\xrightarrow{d}Y$, and $X_n$ and $Y_n$
are independent for each $n$. Show that $(X_n,Y_n)$ converges in
distribution to a random $(k+l)$-vector.
\par\vspace{0.5em}\noindent\textbf{Sol.}\par
Let $\widetilde Y\overset d=Y$ be independent of $X$, constructed
on a product probability space if necessary. At continuity points
$x,y$ of the two marginal CDFs, independence gives
\[
F_{(X_n,Y_n)}(x,y)=F_{X_n}(x)F_{Y_n}(y)
\longrightarrow F_X(x)F_Y(y)=F_{(X,\widetilde Y)}(x,y).
\]
To justify the joint convergence directly, take $s\in\R^k$ and
$t\in\R^l$. The joint characteristic function factorizes as
\[
\begin{aligned}
\E e^{i(s^\top X_n+t^\top Y_n)}
&=\E e^{is^\top X_n}\,\E e^{it^\top Y_n}\\
&\longrightarrow
\E e^{is^\top X}\,\E e^{it^\top\widetilde Y}.
\end{aligned}
\]
The limit is the characteristic function of $(X,\widetilde Y)$
and is continuous at the origin. The multivariate L\'evy continuity
theorem therefore gives
\[
\boxed{(X_n,Y_n)\xrightarrow d(X,\widetilde Y).}
\]
The original $X$ and $Y$ need not themselves be independent;
the conclusion identifies the limit as the product of their laws.

\subsubsection*{Exercise 1.140}
Let $X_n$ be a random variable having the $N(\mu_n,\sigma_n^2)$
distribution, $n=1,2,\ldots$, and $X$ be a random variable having
the $N(\mu,\sigma^2)$ distribution. Show that
$X_n\xrightarrow{d}X$ if and only if $\mu_n\to\mu$ and
$\sigma_n\to\sigma$.

\par\vspace{0.5em}\noindent\textbf{Sol.}\par
Write $\phi_n,\phi$ for the characteristic functions of $X_n,X$.
They are
\[
\phi_n(t)=\exp(i\mu_nt-\sigma_n^2t^2/2),\qquad
\phi(t)=\exp(i\mu t-\sigma^2t^2/2).
\]
Suppose first that $X_n\xrightarrow dX$. Then $\phi_n(t)\to\phi(t)$
for every real $t$. Taking absolute values at $t=1$ gives
\[
e^{-\sigma_n^2/2}\longrightarrow e^{-\sigma^2/2}.
\]
Taking logarithms shows $\sigma_n^2\to\sigma^2$, hence
$\sigma_n\to\sigma$ since standard deviations are nonnegative.
Moreover, for every real $t$,
\[
e^{i\mu_nt}=\phi_n(t)e^{\sigma_n^2t^2/2}
\longrightarrow\phi(t)e^{\sigma^2t^2/2}=e^{i\mu t}.
\]
These are the characteristic functions of the point masses
$\delta_{\mu_n}$ and $\delta_\mu$. By L\'evy's continuity theorem,
$\delta_{\mu_n}\xrightarrow w\delta_\mu$. For every $\epsilon>0$,
the Portmanteau theorem applied to the open interval
$(\mu-\epsilon,\mu+\epsilon)$ gives
\[
\liminf_n\mathbf1_{\{|\mu_n-\mu|<\epsilon\}}\ge1.
\]
Thus $|\mu_n-\mu|<\epsilon$ eventually, proving $\mu_n\to\mu$.
It is important to use all real $t$ here, rather than to take a
single-valued logarithm of a complex phase.

Conversely, if $\mu_n\to\mu$ and $\sigma_n\to\sigma$, the displayed
normal characteristic functions converge pointwise to $\phi$.
Since $\phi$ is continuous at zero, L\'evy's continuity theorem
implies $X_n\xrightarrow dX$.
\vspace{2em}

\subsubsection*{Exercise 1.137 (optional)}
Let $\{X_n\}$ and $\{Y_n\}$ be two sequences of random variables.
Suppose that $X_n\xrightarrow{d}X$ and that
\[
P_{Y_n\mid X_n=x_n}\xrightarrow{w}P_Y
\]
almost surely for every sequence of numbers $\{x_n\}$, where $X$
and $Y$ are independent random variables. Show that
$X_n+Y_n\xrightarrow{d}X+Y$.

\par\vspace{0.5em}\noindent\textbf{Sol.}\par
Use the versions of the conditional distributions appearing in the
hypothesis and write $K_n(x,dy)=P_{Y_n\mid X_n=x}(dy)$.
For a fixed real $t$, let
\[
h_n(x)=\int_{\R}e^{ity}K_n(x,dy),\qquad
\phi_Y(t)=\E e^{itY}.
\]
The hypothesis, together with weak convergence against bounded
continuous functions, implies
\[
h_n(x_n)\longrightarrow\phi_Y(t)
\quad\text{for every sequence }(x_n)
\]
on the sets where the stated conditional-law assumption holds.
Consequently,
\[
\operatorname*{ess\,sup}_{x\ [P_{X_n}]}
|h_n(x)-\phi_Y(t)|\longrightarrow0.
\]
Indeed, otherwise there would be $\epsilon>0$, a subsequence
$n_j$, and points $x_{n_j}$ in the admissible full-measure sets
such that $|h_{n_j}(x_{n_j})-\phi_Y(t)|>\epsilon$.
Completing these points to a sequence contradicts the hypothesis.

By conditioning on $X_n$,
\[
\begin{aligned}
\E e^{it(X_n+Y_n)}
&=\E[e^{itX_n}h_n(X_n)]\\
&=\phi_Y(t)\E e^{itX_n}+o(1)\\
&\longrightarrow\phi_Y(t)\phi_X(t).
\end{aligned}
\]
Since $X$ and $Y$ are independent, the last expression is the
characteristic function of $X+Y$. L\'evy's continuity theorem gives
\[
{X_n+Y_n\xrightarrow dX+Y.}
\]
Here the condition for every sequence of conditioning values is
essential to this argument; mere pointwise convergence for each
fixed conditioning value would not justify substituting $X_n$.
\vspace{2em}

\subsubsection*{Exercise 1.152 (optional)}
Let $T_n=\sum_{i=1}^nX_i$, where the $X_n$'s are independent random
variables satisfying
\[
\P(X_n=n^\theta)=\P(X_n=-n^\theta)=\frac12,
\]
and $\theta>0$ is a constant. Show that
\begin{enumerate}[label=(\alph*),leftmargin=*]
\item when $\theta<0.5$, $T_n/n\xrightarrow{\mathrm{a.s.}}0$;
\item when $\theta\ge1$, $T_n/n\xrightarrow{p}0$ does not hold.
\end{enumerate}

\par\vspace{0.5em}\noindent\textbf{Sol.}\par
Write $X_i=i^\theta\varepsilon_i$, where the independent signs
$\varepsilon_i$ take the values $1,-1$ with equal probability.
Then $\E X_i=0$ and $\operatorname{Var}(X_i)=i^{2\theta}$.
\begin{enumerate}[label=(\alph*),leftmargin=*]
\item If $\theta<1/2$, then
\[
\sum_{i=1}^\infty\operatorname{Var}(X_i/i)
=\sum_{i=1}^\infty i^{2\theta-2}<\infty.
\]
By Kolmogorov's convergence theorem for independent centered variables,
$\sum_iX_i/i$ converges almost surely. Kronecker's lemma then gives
\[
\boxed{\frac{T_n}{n}=\frac1n\sum_{i=1}^nX_i
\xrightarrow{\mathrm{a.s.}}0.}
\]
Equivalently, this is Kolmogorov's sufficient condition for the
strong law of large numbers.

\item Suppose, to the contrary, that $T_n/n\xrightarrow p0$.
For $n\ge2$,
\[
\frac{X_n}{n}
=\frac{T_n}{n}-\frac{n-1}{n}\frac{T_{n-1}}{n-1}
\xrightarrow p0.
\]
But if $\theta\ge1$, then
\[
\left|\frac{X_n}{n}\right|=n^{\theta-1}\ge1
\qquad\text{almost surely},
\]
so $\P(|X_n/n|>1/2)=1$ for every $n$, a contradiction.
Thus $T_n/n\xrightarrow p0$ does not hold.
\end{enumerate}
\vspace{2em}

\subsection*{Problem 2. Asymptotic distributions and estimation}
Bickel and Doksum, \textit{Mathematical Statistics: Basic Ideas and
Selected Topics}, Volume I, second edition, Chapter 5, Section 5.6,
Problems for Section 5.3: Exercises 13 and 33 (required).

\subsubsection*{Exercise 5.3.13}
Let $S_n$ have a $\chi_n^2$ distribution.
\begin{enumerate}[label=(\alph*),leftmargin=*]
\item Show that if $n$ is large, $\sqrt{S_n}-\sqrt n$ has
approximately a $N\left(0,\tfrac12\right)$ distribution. This is known as
\emph{Fisher's approximation}.
\item From (a) deduce the approximation
\[
\P[S_n\le x]\approx\Phi(\sqrt{2x}-\sqrt{2n}).
\]
\item Compare the approximation of (b) with a central limit
approximation
\[
\P[S_n\le x]\approx\Phi\!\left(\frac{x-n}{\sqrt{2n}}\right)
\]
and the exact values of $\P[S_n\le x]$ from the $\chi^2$ table for
$x=x_{0.90}$, $x=x_{0.99}$, and $n=5,10,25$. Here $x_q$ denotes
the $q$th quantile of the $\chi^2$ distribution.
\end{enumerate}

\par\vspace{0.5em}\noindent\textbf{Sol.}\par
\begin{enumerate}[label=(\alph*),leftmargin=*]
\item We complete the characteristic-function argument using the
density of $S_n$. For real $t$,
\[
\phi_{\sqrt{S_n}-\sqrt n}(t)
=\int_0^\infty e^{it(\sqrt x-\sqrt n)}
\frac{x^{n/2-1}e^{-x/2}}{2^{n/2}\Gamma(n/2)}\,dx.
\]
Under the change of variables $u=\sqrt x-\sqrt n$,
we have $x=(u+\sqrt n)^2$ and $dx=2(u+\sqrt n)\,du$.
Hence
\[
\phi_{\sqrt{S_n}-\sqrt n}(t)
=\int_{\R}e^{itu}g_n(u)\,du,
\]
where
\[
g_n(u)=\mathbf1_{\{u>-\sqrt n\}}
\frac{(u+\sqrt n)^{n-1}e^{-(u+\sqrt n)^2/2}}
{2^{n/2-1}\Gamma(n/2)}.
\]
In particular, the power after the substitution is $n-1$.
For a fixed $u$, eventually $u>-\sqrt n$, and
\[
g_n(u)=C_n\exp\left\{
(n-1)\log(1+u/\sqrt n)-u\sqrt n-u^2/2
\right\},
\]
with
\[
C_n=\frac{n^{(n-1)/2}e^{-n/2}}
{2^{n/2-1}\Gamma(n/2)}.
\]
Stirling's formula,
\[
\Gamma(n/2)\sim\sqrt{2\pi}(n/2)^{n/2-1/2}e^{-n/2},
\]
gives $C_n\to1/\sqrt\pi$. Also, Taylor expansion at zero yields
\[
\begin{aligned}
&(n-1)\log(1+u/\sqrt n)-u\sqrt n-u^2/2\\
&\qquad=(n-1)\left(\frac u{\sqrt n}-\frac{u^2}{2n}
+O(n^{-3/2})\right)-u\sqrt n-u^2/2
\longrightarrow -u^2.
\end{aligned}
\]
Thus $g_n(u)\to g(u)=\pi^{-1/2}e^{-u^2}$ for every fixed $u$.
Both $g_n$ and $g$ are probability densities, so Scheff\'e's lemma gives
\[
\int_{\R}|g_n(u)-g(u)|\,du\longrightarrow0.
\]
Consequently,
\[
\left|\int e^{itu}g_n(u)\,du-\int e^{itu}g(u)\,du\right|
\le\int|g_n-g|\longrightarrow0.
\]
The density $g$ is that of $N(0,1/2)$, with characteristic function
$e^{-t^2/4}$. Therefore
\[
\boxed{\sqrt{S_n}-\sqrt n\xrightarrow dN(0,1/2).}
\]
This justifies convergence of the integrals without an unsupported
interchange of a pointwise limit and integration.

An alternative short proof is to write $S_n$ in distribution as
$\sum_{i=1}^nG_i^2$ for i.i.d.\ standard normal $G_i$. The CLT gives
$(S_n-n)/\sqrt{2n}\xrightarrow dN(0,1)$, while $S_n/n\xrightarrow p1$.
Therefore
\[
\sqrt{S_n}-\sqrt n
=\frac{(S_n-n)/\sqrt n}{\sqrt{S_n/n}+1}
\xrightarrow dN(0,1/2)
\]
by Slutsky's theorem.

\item Standardizing the approximation in (a), for $x\ge0$,
\[
\begin{aligned}
\P(S_n\le x)
&=\P\bigl(\sqrt2(\sqrt{S_n}-\sqrt n)
\le\sqrt{2x}-\sqrt{2n}\bigr)\\
&\approx\Phi(\sqrt{2x}-\sqrt{2n}).
\end{aligned}
\]
For $x<0$, the exact CDF is zero. The displayed approximation is
intended for the central and relevant tail regions as $n$ becomes large.

\item Let $x_q$ be the $q$th quantile of $\chi_n^2$.
The exact CDF at this value is $q$. The following table compares
\[
F_{\mathrm F}(x)=\Phi(\sqrt{2x}-\sqrt{2n}),\qquad
F_{\mathrm C}(x)=\Phi\left(\frac{x-n}{\sqrt{2n}}\right).
\]
Values are rounded to six decimal places; the last two columns are
absolute errors relative to $q$.
\begin{center}
\small
\setlength{\tabcolsep}{4pt}
\begin{tabular}{rrrrrrr}
\hline
$n$ & $q$ & $x_q$ & $F_{\mathrm F}(x_q)$ & $F_{\mathrm C}(x_q)$
& $|F_{\mathrm F}-q|$ & $|F_{\mathrm C}-q|$\\
\hline
5 & 0.90 & 9.236357 & 0.871961 & 0.909821 & 0.028039 & 0.009821\\
5 & 0.99 & 15.086272 & 0.990115 & 0.999288 & 0.000115 & 0.009288\\
10 & 0.90 & 15.987179 & 0.881487 & 0.909678 & 0.018513 & 0.009678\\
10 & 0.99 & 23.209251 & 0.990383 & 0.998430 & 0.000383 & 0.008430\\
25 & 0.90 & 34.381587 & 0.889012 & 0.907705 & 0.010988 & 0.007705\\
25 & 0.99 & 44.314105 & 0.990440 & 0.996847 & 0.000440 & 0.006847\\
\hline
\end{tabular}
\end{center}
At the $0.90$ quantile, the central limit approximation is more
accurate for all three sample sizes. At the $0.99$ quantile,
Fisher's approximation is substantially more accurate.
Thus neither approximation is uniformly better at all the
quantiles considered.
\end{enumerate}
\vspace{2em}

\subsubsection*{Exercise 5.3.33}
Let $X_{ij}$, $i=1,\ldots,p$, $j=1,\ldots,k$, be independent with
$X_{ij}\sim N(\mu_i,\sigma^2)$.
\begin{enumerate}[label=(\alph*),leftmargin=*]
\item Show that the MLEs of $\mu_i$ and $\sigma^2$ are
\[
\widehat\mu_i=\frac1k\sum_{j=1}^kX_{ij},
\qquad
\widehat\sigma^2=\frac1{kp}\sum_{i=1}^p\sum_{j=1}^k
(X_{ij}-\widehat\mu_i)^2.
\]
\item Show that if $k$ is fixed and $p\to\infty$, then
\[
\widehat\sigma^2\xrightarrow{p}\frac{k-1}{k}\sigma^2.
\]
That is, the MLE $\widehat\sigma^2$ is not consistent
(Neyman and Scott, 1948).
\item Give a consistent estimate of $\sigma^2$.
\end{enumerate}

\medskip
\noindent\textbf{Correction to part (b):}
The inconsistent estimator in the final sentence is the variance MLE
$\widehat{\sigma}^{\,2}$. Use $k\ge2$ observations per group, with $k$
fixed as $p\to\infty$.

\par\vspace{0.5em}\noindent\textbf{Sol.}\par
Throughout, take $k\ge2$ and $\sigma^2>0$. Write $s=\sigma^2$.
\begin{enumerate}[label=(\alph*),leftmargin=*]
\item The likelihood and log likelihood are
\[
\begin{aligned}
L(\mu_1,\ldots,\mu_p,s)
&=(2\pi s)^{-kp/2}
\exp\left[-\frac1{2s}\sum_{i=1}^p\sum_{j=1}^k
(X_{ij}-\mu_i)^2\right],\\
\ell(\mu_1,\ldots,\mu_p,s)
&=-\frac{kp}{2}\log(2\pi s)
-\frac1{2s}\sum_{i=1}^p\sum_{j=1}^k(X_{ij}-\mu_i)^2.
\end{aligned}
\]
Let $\overline X_i=k^{-1}\sum_{j=1}^kX_{ij}$.
For every $s>0$, the identity
\[
\sum_{j=1}^k(X_{ij}-\mu_i)^2
=\sum_{j=1}^k(X_{ij}-\overline X_i)^2
+k(\overline X_i-\mu_i)^2
\]
shows that the unique maximizing mean is
$\widehat\mu_i=\overline X_i$.
Set
\[
Q=\sum_{i=1}^p\sum_{j=1}^k(X_{ij}-\overline X_i)^2.
\]
After substitution, the derivative of the profile log likelihood is
\[
\ell_{\mathrm{prof}}'(s)
=-\frac{kp}{2s}+\frac{Q}{2s^2}
=\frac{Q-kps}{2s^2}.
\]
It is positive for $s<Q/(kp)$ and negative for $s>Q/(kp)$.
Since $Q>0$ almost surely for $k\ge2$, this proves the global maximum:
\[
\boxed{\widehat\mu_i=\frac1k\sum_{j=1}^kX_{ij},\qquad
\widehat\sigma^2=\frac Q{kp}.}
\]

\item For each group, let
\[
Q_i=\sum_{j=1}^k(X_{ij}-\overline X_i)^2.
\]
Writing $X_{ij}=\mu_i+\sigma Z_{ij}$ removes $\mu_i$, where all the
$Z_{ij}$ are independent standard normal variables. The vector of
centered observations within a group is the orthogonal projection
onto the $(k-1)$-dimensional subspace perpendicular to $(1,\ldots,1)$.
Rotational invariance of the standard normal distribution therefore gives
\[
\frac{Q_i}{\sigma^2}\sim\chi_{k-1}^2.
\]
The $Q_i$ are independent and identically distributed, even though
the group means $\mu_i$ need not be equal. In particular,
\[
\E Q_i=(k-1)\sigma^2,\qquad
\operatorname{Var}(Q_i)=2(k-1)\sigma^4.
\]
By the strong law of large numbers, as $p\to\infty$ with $k$ fixed,
\[
\widehat\sigma^2
=\frac1k\frac1p\sum_{i=1}^pQ_i
\xrightarrow{\mathrm{a.s.}}\frac{k-1}{k}\sigma^2.
\]
This implies the requested convergence in probability. Because the
limit is not $\sigma^2$, the variance MLE is inconsistent.
Its exact bias is $-\sigma^2/k$, which does not disappear when
only the number of groups grows.

\item Correct the degrees of freedom by defining
\[
{\widetilde\sigma^2
=\frac{k}{k-1}\widehat\sigma^2
=\frac1{p(k-1)}\sum_{i=1}^p\sum_{j=1}^k
(X_{ij}-\overline X_i)^2.}
\]
Then $\E\widetilde\sigma^2=\sigma^2$, and
\[
\operatorname{Var}(\widetilde\sigma^2)
=\frac{2\sigma^4}{p(k-1)}\longrightarrow0.
\]
Chebyshev's inequality proves
$\widetilde\sigma^2\xrightarrow p\sigma^2$; in fact, the same
strong-law argument proves almost sure convergence.
\end{enumerate}
\vspace{2em}

\end{document}
